\documentclass[10pt,a4paper]{amsart}
\usepackage[margin=19mm]{geometry}
\usepackage{amsmath,amssymb,mathtools}
\usepackage[scr=rsfs,cal=euler]{mathalfa}
\usepackage{xfrac}
\usepackage[unicode,hidelinks,bookmarks=false]{hyperref}
\newcommand{\ie}{i.e.\ }
\newcommand{\cf}{cf.\ }
\newcommand{\resp}{respectively\ }
\newcommand{\Subsection}{\subsection}
\providecommand{\nobreakdash}{\nobreak\hbox{-}}
\setlength{\emergencystretch}{2em}
\multlinegap0pt
\usepackage{amssymb}
\usepackage{dsfont}
\newtheorem{lemma}{Lemma}
\def\Zv{\mathscr{Z}}
\def\be{\begin{equation}}
\newcommand{\bfv}{{\boldsymbol v}}
\newcommand{\cI}{\mathcal{I}}
	\def\ee{\end{equation}}
\title{Remarks on Vafa-Witten theory, \\ or how to sit on a wall}
\author{Sergei Alexandrov, Aradhita Chattopadhyaya}
\date{Source: arXiv:2608.25007v1 (CC BY 4.0)}
\begin{document}
\maketitle

\noindent\emph{Source and licence.} Remarks on Vafa-Witten theory, \\ or how to sit on a wall. Version arXiv:2608.25007v1 (CC BY 4.0), distributed by arXiv under the Creative Commons Attribution 4.0 International licence, http://creativecommons.org/licenses/by/4.0/, as recorded in the arXiv licence metadata for this exact version and shown on the abstract page https://arxiv.org/abs/2608.25007v1. Full licence notice: \texttt{licenses/arxiv-2608.25007-CC-BY-4.0.txt}.\par\smallskip

\noindent\textit{Editorial scope.} Counted unit: the article's lemma lemma1 (source0.tex lines 1624--1635). The article's complete original proof of it is delivered with it (source0.tex lines 1636--1658, one \texttt{proof} environment of 23 lines). No mathematics is written, repaired or re-derived: the statement and the proof are the article's own lines, in the article's own notation, and no retained line is substituted or re-flowed.  No further source-local context is needed: every cross-reference of every retained line resolves inside the two delivered spans.  Nothing was omitted from any retained span, so the package carries no omission record.  No retained line contains a citation, so the wrapper carries no bibliography and no bibliography entry.  No retained line contains a figure environment, an image-inclusion command or drawing code, so no image is delivered and the package declares no asset dependency.  Editorial formatting only, in addition: an AMS \texttt{amsart} preamble that restates the article's own declarations and macros used by the retained lines, namely the environments \texttt{lemma} on separate counters, together with the standard packages those lines need.  The retained environments print in this document as Lemma 1 in order of appearance, whereas the article prints the same items with the numbering of its own full document.  The complete native source of the article as distributed in the arXiv e-print for this exact version is bundled unchanged as \texttt{sources/208535/source0.tex}.
\begin{lemma}
	\label{lemma1}
	Let $\cI\subset \Zv_{n-1}$ and $i\notin \cI$. We denote $i_-=\max\{j\in \cI\cup\{0\}:\ j<i\}$
	and $i_+=\min\{j\in \cI\cup\{n\}:\ j>i\}$. Besides, $\bfv_{i\perp \cI}$ denotes the projection of $\bfv_i$ 
	orthogonal to the subspace spanned by $\{\bfv_j\}_{j\in\cI}$.	
	Then in the collinear case where $p_i=r_ip_0$, one has 
	\be
	\bfv_{\ell\perp \cI}=\lambda \sum_{i=i_-+1}^\ell\sum_{j=\ell+1}^{i_+}\bfv_{ij},
	\qquad 
	\lambda=\frac{r}{\sum_{l=i_-+1}^{i_+}r_l}\, .
	\ee
\end{lemma}

\begin{proof}
	It is straightforward to calculate that 
	\be
	\begin{split}
		\bfv_{\ell\perp k} =&\, \frac{r}{\sum_{l=1}^k r_l}\sum_{i=1}^\ell \sum_{j=\ell+1}^k \bfv_{ij},
		\qquad\quad \
		k>\ell,
		\\  
		\bfv_{\ell\perp k} =&\, \frac{r}{\sum_{l=k+1}^n r_l}\sum_{i=k+1}^\ell \sum_{j=\ell+1}^n \bfv_{ij},
		\qquad 
		k<\ell. 
	\end{split}
	\label{iperpj}
	\ee 
	Since $\bfv_\ell\ast\bfv_{ij}=0$ for $\ell\ne i,j$, 
	this result ensures that $\bfv_{\ell\perp i}\ast \bfv_j=0$ both for $\ell<i<j$ and for $j<i<\ell$, and thus already proves the claim
	in the cases of $i_-=0$ and $i_+=n$. In other cases, we compute 
	\be
		\bfv_{\ell\perp \{i_-,i_+\}} = \frac{r}{\sum_{l=i_-+1}^{i_+}r_l}\sum_{i=i_-+1}^\ell\sum_{j=\ell+1}^{i_+}\bfv_{ij}.
	\ee 
	Since this result satisfies $\bfv_{\ell\perp \{i_-,i_+\}}\ast \bfv_j=0$ for $j<i_-$ and $j>i_+$, we conclude that
	$\bfv_{\ell\perp \cI}=\bfv_{\ell\perp \{i_-,i_+\}}$, which completes the proof.
\end{proof}

\end{document}
